\section{The automorphic projection and its pointwise trace}
\label{ak-section}

We supply the kernel construction required in
Section~\ref{amb-section}.  The diagonal identity is proved
pointwise.  Equality of measurable kernels almost everywhere
on a product would not suffice for this purpose.
The only additional arithmetic theorem used here is the
classical finiteness theorem for automorphic forms of fixed
level, compact type and infinitesimal character; its precise
use is stated below.

\subsection{A matrix coefficient gives the required projection}

All Hilbert space inner products in this section are linear
in the second variable.  Let \(G\) be a unimodular locally
compact group, let \(\sigma\) be an irreducible unitary
representation on a Hilbert space \(H\), and let \(w\in H\)
have norm one.  Suppose its formal degree is \(d>0\).
Polarization of the formal-degree identity gives
\[
 d\int_G\overline{\langle a,\sigma(g)b\rangle}
                 \langle c,\sigma(g)v\rangle\,dg
                       =\langle c,a\rangle\langle b,v\rangle .
 \tag{AK1}
\]
Every integral in (AK1) is absolutely convergent by
Cauchy--Schwarz.  Put
\[
 a_w(g)=\langle w,\sigma(g)w\rangle,\qquad
                   \phi(g)=d\,\overline{a_w(g)}.
 \tag{AK2}
\]
Assume \(\phi\in L^1(G)\).

\begin{proposition}\label{ak-projection}
For every strongly continuous unitary representation \(\pi\)
on a Hilbert space \(\mathcal H\), the bounded operator
\[
                  P=\int_G\phi(g)\pi(g)\,dg
\]
is an orthogonal projection.  Moreover
\[
                  P\pi(x)P=a_w(x)P \qquad(x\in G).
 \tag{AK3}
\]
Evaluation at \(w\) is a linear bijection
\[
 \operatorname{Hom}^{\rm cont}_G(H,\mathcal H)
                       \longrightarrow\operatorname{ran}P,
                  \qquad T\longmapsto Tw .
 \tag{AK4}
\]
The intertwiner corresponding to a vector \(v\) has operator
norm \(\|v\|\).
\end{proposition}
\begin{proof}
The operator integral exists on each vector, by strong
continuity and the bound \(\|\pi(g)v\|=\|v\|\).
Its norm is at most \(\|\phi\|_1\).
Unitarity gives
\(\phi(g^{-1})=\overline{\phi(g)}\), so \(P=P^*\).
Apply (AK1) with \(a=b=w\), \(c=\sigma(y)w\), and
\(v=\sigma(x)w\).  The result is
\[
 \int_G\phi(g)\phi(x^{-1}g^{-1}y)\,dg
                                      =a_w(x)\phi(y).
 \tag{AK5}
\]
Indeed the second factor in the integrand is
\(d\langle\sigma(y)w,\sigma(g)\sigma(x)w\rangle\).
For \(x=1\), this proves \(\phi*\phi=\phi\).
Integrating (AK5) in \(\pi\), using Fubini for the
\(L^1\) functions, proves both \(P^2=P\) and (AK3).

If \(v\in\operatorname{ran}P\) has norm one, (AK3) gives
\(\langle v,\pi(g)v\rangle=a_w(g)\).  Therefore the map
\[
       \sum_i c_i\sigma(g_i)w\longmapsto
                          \sum_i c_i\pi(g_i)v
\]
preserves the squared norm: expand it as a finite sum
of the coefficients at \(g_i^{-1}g_j\).
It is well defined and extends to an isometry \(T_v\)
on \(H\), since the orbit of the nonzero vector \(w\)
is dense by irreducibility.  The same formula proves
equivariance.  Scaling treats every \(v\), including zero.
Conversely, any continuous intertwiner satisfies
\(P(Tw)=T\sigma(\phi)w=Tw\), by (AK1).
Density of the orbit of \(w\) proves uniqueness.
The isometry construction gives the asserted operator norm.
\end{proof}

We apply this to the product of the archimedean
representations in Section~\ref{amb-section}.
Write \(d_\infty=\prod_vd_v\) and
\(w=\bigotimes_v w_v\).  The finite tensor product has
formal degree \(d_\infty\): multiply (AK1) for elementary
tensors and use their density.  Let \(E_f=R(e)\) on
\(L^2(\mathcal X)\).  The finite convolution identities
in \eqref{amb-idempotent} show that \(E_f\) is an
orthogonal projection.  It commutes with every
archimedean right translation.  Hence
\[
 Q=E_fP=R(\Phi),\qquad
 \Phi(g_\infty,g_f)=d_\infty
       \overline{\langle w,\sigma(g_\infty)w\rangle}\,e(g_f)
 \tag{AK6}
\]
is also an orthogonal projection.  Formula (AK4) restricts
to a linear bijection between its image and the intertwiner
space in \eqref{amb-D-definition}: the condition on an
intertwiner is \(E_fT=T\).  Thus the rank of \(Q\), once
proved finite, is exactly the required \(D\).

\subsection{Two-point counting and convergence of the series}

Retain \(n,f,q,h\) and \(\tau_x\) from
Section~\ref{amb-section}, with \(n\ge3\) and \(q\le n/2\),
as in its large-\(n\) application.  We first record two extensions
of its lattice argument.

\begin{lemma}\label{ak-count}
If \(T:\Lambda_y\to\Lambda_x\) is a nonzero
\(\mathcal O_E\)-linear map, then
\[
              \left(\prod_v\|T_v\|_{\rm op}\right)^{1/f}
                                      \ge\lambda_x\lambda_y.
 \tag{AK7}
\]
Consequently the number of rational \(\gamma\) for which
\((x^{-1}\gamma y)_f\in K_f\) and
\(t((x^{-1}\gamma y)_v)\le r_v\), with all \(r_v\ge0\),
is at most
\[
 \left(\frac{3\sqrt n}{\tau_x\tau_y}\right)^{2n^2f}
                               \exp\left(2n^2\sum_v r_v\right).
 \tag{AK8}
\]
The function \(x\mapsto\tau_x\) is positive and continuous.
\end{lemma}
\begin{proof}
In the proof of Lemma~\ref{amb-map-separation}, use
the quotient of \(\Lambda_y\) by the kernel, and the
image in \(\Lambda_x\).  The quotient size is at most
\(\lambda_y^{-1}\), and the image size is at least
\(\lambda_x\).  The same covolume inequality proves (AK7).
The difference of two rational matrices counted in (AK8)
is such a map.  Lemma~\ref{amb-packing}, in the
\(2n^2\)-dimensional real matrix spaces, proves (AK8)
exactly as it proved \eqref{amb-ball-bound}.

The finite point lattice is locally constant modulo \(K_f\).
At infinity, a change of positive metric by maps whose
norms and inverse norms are at most \(e^{r_v}\)
changes every sublattice size by a factor between
\(e^{-\sum_vr_v/f}\) and \(e^{\sum_vr_v/f}\).
Taking infima preserves these bounds.  Taking the minimum
with \(1/2\) preserves them as well.  Since the maps tend
to the identity when the point tends to a fixed point,
these inequalities prove continuity.  Positivity was proved
in Lemma~\ref{amb-size-duality}.
\end{proof}

For \(q>0\), assume
\[
                         H=\frac{h}{8\sqrt{2q}}>2n^2.
 \tag{AK9}
\]
This is much weaker than \eqref{amb-parameters}.
The proved compact-type bound implies, at each place,
\[
 \bigl|\langle w_v,\sigma_v(g)w_v\rangle\bigr|
 \le \exp\!\left(-\frac h8
                \min\{t(g)^2,t(g)/\sqrt{2q}\}\right)
                         \le e^H e^{-Ht(g)}.
 \tag{AK10}
\]
For \(t\le1\), the last right side is at least one;
for \(t\ge1\), the minimum in the exponent is
\(t/\sqrt{2q}\).  This proves the second inequality.

The coefficient is integrable.  Here is the growth
calculation needed to justify that assertion.  In the
Cartan coordinates \(a_1\ge\cdots\ge a_q\ge0\) on
\(U(n-q,q)\), one has \(t^2=2\sum_i a_i^2\).
The positive restricted roots have multiplicities
\[
 e_i\pm e_j:2,\qquad e_i:2(n-2q),\qquad 2e_i:1.
\]
The Jacobian of the polar map, up to a fixed positive
constant depending on the Haar normalization, is
the product of \(\sinh(\alpha(a))^{m_\alpha}\).
This follows by differentiating the polar map on each
restricted-root subspace: its two opposite directions
give the factor \(\sinh(\alpha(a))\).
Bounding each such factor by its exponential gives
\[
 J(a)\le C\exp\left(2\sum_i(n-2i+1)a_i\right)
                                \le C e^{\sqrt2 n\sqrt q\,t}.
\]
Since \(q\le n\), condition (AK9) implies
\(H>\sqrt2 n\sqrt q\).  Integrating (AK10) over the
positive chamber, and then over the two compact factors,
is finite; the remaining radial volume is at most a
constant times \(t^{q-1}\,dt\).  When \(q=0\) the
archimedean group is compact, so integrability is immediate.
The finite function \(e\) is compactly supported.
Thus \(\Phi\in L^1(G_\infty\times G_f)\), as required
for (AK6).

\begin{theorem}\label{ak-kernel}
Under (AK9), or in the compact case \(q=0\), the series
\[
                    K(x,y)=\sum_{\gamma\in\Gamma}
                                        \Phi(x^{-1}\gamma y)
 \tag{AK11}
\]
converges absolutely and locally uniformly.
It defines a continuous kernel on
\(\mathcal X\times\mathcal X\), and for some finite
constant \(C_\Phi\),
\[
 \sum_\gamma|\Phi(x^{-1}\gamma y)|
                  \le C_\Phi(\tau_x\tau_y)^{-2n^2f}.
 \tag{AK12}
\]
It represents \(Q\) on continuous compactly supported
test functions.
\end{theorem}
\begin{proof}
Positivity of \(e\) gives \(|e(g_f)|\le e(1)\), by its
two-by-two positive matrix, as in
Lemma~\ref{amb-positive-type}.
For \(q>0\), separate the series into the shells
\(j_v\le t((x^{-1}\gamma y)_v)<j_v+1\),
where all \(j_v\) are nonnegative integers.
Equations (AK8)--(AK10) bound its absolute value by
\[
 m_0\left(\frac{3\sqrt n}{\tau_x\tau_y}\right)^{2n^2f}
 e^{f(H+2n^2)}
 \sum_{(j_v)\in\mathbb Z_{\ge0}^{\,f}}
                             e^{-(H-2n^2)\sum_vj_v}.
\]
The sum is \((1-e^{-(H-2n^2)})^{-f}\).
This proves (AK12), including \(e=0\).
For \(q=0\), all distances are zero and (AK8)
directly bounds the finite sum.

On a compact set of pairs, \(\tau_x\tau_y\) has a positive
lower bound.  The same geometric sum bounds the tails
uniformly.  The union of the rational elements with all
distances bounded on such a compact set is finite:
their finite components lie in one fixed compact adelic
set and their archimedean matrices are bounded.
After a fixed denominator is cleared, their entries
belong to a fixed lattice in the Minkowski space of \(E\),
whose bounded subsets are finite.  Thus these are tails
of the original series, not only tails of a varying
partition.  Local uniform convergence and continuity
of each term prove continuity of \(K\).
Reindexing \(\Gamma\) proves its invariance in both
quotient variables, and inversion proves
\(K(y,x)=\overline{K(x,y)}\).

For \(F\in C_c(\mathcal X)\), unfolding gives
\[
 \int_{\mathcal X}K(x,y)F(y)\,dy
       =\int_{G_\infty\times G_f}\Phi(g)F(xg)\,dg .
 \tag{AK13}
\]
Absolute convergence follows from
\(\|F\|_\infty\|\Phi\|_1<\infty\).
The latter integral represents the bounded \(L^2\)
operator \(R(\Phi)\), by its defining operator integral
and Fubini against a compactly supported test function.
The left side is continuous locally in \(x\), by
continuity of \(K\) on compact sets.  Thus (AK13)
gives a continuous representative of \(QF\).
\end{proof}

\subsection{A continuous kernel determines its diagonal}

\begin{lemma}\label{ak-diagonal}
Let \(X\) be a locally compact second countable space
with a positive Radon measure of full support.
Let \(Q\) be an orthogonal projection on \(L^2(X)\),
represented on \(C_c(X)\) by a continuous Hermitian
kernel \(K\).  Then every vector of \(\operatorname{ran}Q\)
has a unique continuous representative.  For an orthonormal
basis \((F_i)\) of that range, with these representatives,
\[
 K(x,y)=\sum_iF_i(x)\overline{F_i(y)},\qquad
 K(x,x)=\sum_i|F_i(x)|^2\ge0
 \tag{AK14}
\]
at every point.  In particular
\[
                \int_X K(x,x)\,dx
                    =\dim\operatorname{ran}Q
 \tag{AK15}
\]
as an equality of nonnegative extended real numbers.
\end{lemma}
\begin{proof}
Choose nonnegative \(u_{x,j}\in C_c(X)\), of integral
one, supported in successively smaller neighborhoods
of \(x\).  Such functions exist by local compactness,
full support and regularity of the measure.  The kernel
formula and continuity give
\[
 \langle Qu_{x,j},Qu_{y,l}\rangle
   =\langle u_{x,j},Qu_{y,l}\rangle
                   \longrightarrow K(x,y).
\]
Taking \(y=x\) shows that \(Qu_{x,j}\) is a Cauchy
sequence in \(L^2\).  Let its limit be \(k_x\).
It lies in the closed range of \(Q\), and
\[
                  \langle k_x,k_y\rangle=K(x,y).
 \tag{AK16}
\]
In particular \(x\mapsto k_x\) is continuous in Hilbert
norm, by expanding \(\|k_x-k_y\|^2\).

For \(F\in\operatorname{ran}Q\), define
\(\widetilde F(x)=\langle k_x,F\rangle\).
This is continuous.  It represents \(F\) in \(L^2\).
To verify this last assertion without assigning point
values to an \(L^2\) class, first let \(u\in C_c(X)\).
Pairing with \(v\in C_c(X)\) in the defining limit gives
\(\langle v,k_x\rangle=\overline{(Qv)(x)}\).
Integrating this identity and using \(Q=Q^*\) gives
\[
                Qu=\int_X u(x)k_x\,dx
\]
as a Hilbert space integral; the integrand is continuous
and bounded on its compact support.
Pairing with \(F\) gives
\(\int\overline{u(x)}\widetilde F(x)\,dx
=\langle Qu,F\rangle=\langle u,F\rangle\).
Thus \(\widetilde F=F\) almost everywhere.
Two continuous representatives agreeing almost everywhere
are equal, because a nonempty open set has positive measure.

Expand \(k_x\) and \(k_y\) in an orthonormal basis of
\(\operatorname{ran}Q\).  Parseval and Cauchy--Schwarz
give both formulas in (AK14), with an absolutely convergent
series at each pair of points.  Separability gives a
countable basis.  Tonelli and \(\|F_i\|_2=1\) prove
(AK15), including an infinite-dimensional range.
\end{proof}

\subsection{Finiteness and the exact trace}

We use the following classical arithmetic input.
For a reductive group over a number field, the space
of smooth automorphic forms of fixed compact open
level, fixed finite collection of archimedean compact
types, and annihilated by a fixed finite-codimensional
ideal of the center of the universal enveloping algebra,
is finite dimensional.  Moderate growth is required.
This is Harish--Chandra's finiteness theorem; the adelic
version follows from Borel's finite class decomposition.
The precise statements and the passage from the
classical to the adelic quotient are given in
\cite[Theorems 6.2.3 and 6.3.2]{getz-hahn-2022}.
These classical theorems are imported here.

\begin{theorem}\label{ak-finite-trace}
In the setting of Section~\ref{amb-section}, under (AK9)
or \(q=0\), the operator \(Q=R(\Phi)\) has finite rank
\(D\) as defined in \eqref{amb-D-definition}.
Its continuous kernel satisfies
\[
       K(x,x)=\sum_{i=1}^D|F_i(x)|^2,\qquad
                      \int_{\mathcal X}K(x,x)\,dx=D
 \tag{AK17}
\]
for a continuous orthonormal basis of its image.
\end{theorem}
\begin{proof}
It remains to check the hypotheses of the stated
finiteness theorem; neither finite rank nor an
integrable diagonal is assumed in this check.
By Proposition~\ref{ak-projection}, every vector
of the range is \(Tw\) for a continuous intertwiner
\(T\) satisfying \(E_fT=T\).
It is fixed by \(K_f(g)\), because \(e\) is
bi-invariant under that subgroup.
Its archimedean compact translates lie in the
image under \(T\) of the fixed lowest compact type.
The center of the universal enveloping algebra
acts on that image by the infinitesimal character
of \(\sigma\).

We justify smoothness directly.  The lowest compact
type occurs with multiplicity one in \(\sigma\).
Compact Schur orthogonality supplies a matrix
coefficient measure \(p_w\) on \(K_\infty\)
whose action on \(\sigma\) is projection onto
\(\mathbb Cw\).  Choose \(b\in C_c^\infty(G_\infty)\)
with \(c=\langle w,\sigma(b)w\rangle\ne0\);
an approximate identity supplies such \(b\).
Then
\[
                  a=c^{-1}p_w*b*p_w\in C_c^\infty(G_\infty),
                  \qquad \sigma(a)w=w.
\]
Consequently \(R(a)Tw=Tw\).
Convolution by \(a\), together with normalized Haar
measure on \(K_f(g)\), gives smooth representatives
of these \(L^2\) vectors.  To see this in local
coordinates, its kernel is a locally finite sum of
smooth functions: on compact sets all possible
rational matrices lie in a fixed compact adelic set.
For \(x\) in a fixed compact neighborhood the
other variable is restricted to the image of
that neighborhood times the fixed compact support.
Every derivative of the kernel is bounded there,
so Cauchy--Schwarz justifies differentiation under
the integral.  These representatives agree with
the continuous representatives in
Lemma~\ref{ak-diagonal}.
The compact-type and infinitesimal-character
identities therefore hold as identities of smooth
functions, not only of Hilbert space vectors.

Finally (AK12) and (AK14) give
\[
                 |(Tw)(x)|\le
                  C_\Phi^{1/2}\tau_x^{-2n^2f}\|Tw\|_2 .
 \tag{AK18}
\]
On a fixed finite adelic component, let \(\Lambda_*\)
be its point lattice at archimedean identity.
Changing its metric by \(g_\infty^{-1}\) gives
\[
       \tau_{(g_\infty,x_f)}
             \ge\tau_{(1,x_f)}
                   \left(\prod_v\|g_v\|_{\rm op}\right)^{-1/f}.
\]
The operator norms are at least one, since the
matrices are unitary for an involutory Hermitian form.
Thus (AK18) is polynomial growth in matrix coordinates
on every classical component.  There are finitely many
such components at the fixed level by Borel's class
decomposition.  Apply the classical finiteness theorem
on each component, with the same compact type and
the kernel of the fixed infinitesimal character.
Their finite direct sum contains the range of \(Q\).
This proves finite rank.
The bijection (AK4) identifies the rank with \(D\).
Now (AK17) follows from Lemma~\ref{ak-diagonal}.
\end{proof}

This proves the kernel and trace input in
Proposition~\ref{amb-final-reduction}.  Its nearby
unipotent, small cusp-volume and deep cusp-mass bounds
are separate quantitative statements; the qualitative
finiteness theorem gives none of those bounds.
The construction also explains the role of continuity:
on a nonatomic space an arbitrary measurable kernel
may be changed on the diagonal without changing its
integral operator.  No such change preserves the
continuous kernel constructed above.


% A direct proof of the volume estimate for small least size.
% All labels in this fragment have prefix rv-.
